\section{Discussion}
\label{sec:discussion}

We have described how the error of a case-count estimate of the reproduction number depends on case volume when transmission varies between periods, offspring numbers are overdispersed and only a proportion of infections is reported. For the period-specific value, error falls without limit as reported cases accumulate. For the mean across periods and for one-step prediction, error per period approaches the environmental variance, and the balance scale $\mx=A/v_R$ marks the case volume beyond which each doubling of cases removes at most a quarter of the remaining variance. In weekly influenza hospital admissions over four seasons, two apparent statistics had positive intercepts, and a two-step inversion through a data-matched simulation mapped both to a latent environmental variance of about 0.04. A national-week component accounted for about 40\% to 50\% of the intercept, and reporting completeness under mandatory reporting contributed less than an equivalence margin of 17\% of the intercept, set at a reference reporting volatility of 0.10.

The theory extends a standard result in population ecology to surveillance data. Ecological models split the variance of population growth into a demographic term that falls with population size and an environmental term that does not \citep{engen1998demographic, lande2003stochastic}. The present model retains this partition and adds four elements that matter for surveillance. Transmission is a branching process with negative binomial offspring, so that superspreading enlarges the demographic term \citep{lloydsmith2005superspreading}. Binomial reporting adds a further term that falls with case volume, and reporting volatility adds one that does not (Proposition~\ref{prop:report-noise}). The quantity of interest is an estimate averaged over periods, so the balance scale also sets the optimal division of a budget between cases per period and number of periods (Proposition~\ref{prop:allocation}). That rule has the form of optimal allocation in two-stage sampling \citep{cochran1977sampling}, with the balance scale in place of the ratio of within-unit to between-unit variance.

The results complement the information-theoretic analyses of Parag and colleagues \citep{parag2021improved, parag2022quantifying}. Their metric for a noisy epidemic curve already contains time-varying reporting fractions, which enter through the geometric mean of the reporting probabilities, and a negative binomial extension is treated in their supplementary material \citep{parag2022quantifying}. Their target is the sequence of period-specific values, and in a Poisson renewal model the error variance of each estimate is $\Rbar/m$, so that reporting acts by rescaling the effective case count. Here the target is the mean across periods, and variation shared by all chains of transmission in a period adds a term that does not fall with $m$. The two descriptions agree at small counts and separate near the balance scale. On the apparent scale of $R_A$, the fitted variance was twice the Poisson variance at 22 admissions per week and ten times at 261, and on the latent scale the corresponding values were 28 and 262. Both descriptions are therefore needed: the Poisson analysis ranks data streams by the information they carry about current transmission, and the variance partition here indicates when further cases cease to improve estimates of typical transmission.

The two-step inversion connects the apparent statistics with the latent model. The simulation first established that neither statistic estimates $v_R/\Rbar^2$ directly; the intercept of $G$ was 2.7 times this quantity and that of $R_A$ about half of it, for reasons that follow from how each statistic aggregates growth over a week. After calibration, the two statistics implied latent variances of 0.043 and 0.040 (ratio 0.93; 95\% CI: 0.78, 1.18). This consistency is expected if both statistics respond to the same latent variation and the simulation captures how each transforms it. As a test of the calibration it has limited power: across the simulated settings it failed only when a synchronous shock to reported counts replaced the shock to transmission, and in that setting the observed lag-1 autocorrelation of $G$ lay above 99\% of simulated panels. The simulated panels further reproduced the intraclass correlation, the national-week share and the lag-1 autocorrelation of the data. Within the stated uncertainties, the calibration is therefore consistent with the empirical intercepts, under three conditions. The implied variance depends on the true generation interval. Across the 95\% interval of 2.9 to 4.3 days for the serial interval of influenza \citep{cowling2009estimation}, it ranged from 0.026 to 0.063, so the inversion identifies $v_R$ conditionally on that interval. Weekly shocks were taken as independent. The observed autocorrelation lay near the centre of the simulated distribution under independence, and below 97\% of simulated panels with a first-order coefficient of 0.25 for $R_A$; a coefficient of 0.5 would have lowered the implied variance to about 0.026. Finally, the calibrated theoretical $b/a$ lies below both empirical point estimates, which fall at the 91st and 93rd percentiles of the simulated sampling distribution. Overdispersion in weekly admissions can account for this excess. With a variance-to-mean ratio of 1.5, the simulated $b/a$ rose to 68.0 for $R_A$ and 40.2 for $G$, above both empirical values, while the implied variance stayed at 0.040 and 0.044. A true generation interval of 2.9 days also reproduced the empirical $b/a$, with an implied variance of 0.026. The excess can therefore come from overdispersion smaller than a variance-to-mean ratio of 1.5 or from a shorter generation interval. The two explanations imply variances of about 0.040 and 0.026. The autocorrelation of the data argues against the shorter interval but does not exclude it. The data cannot exclude a mixture of transmission and reporting shocks (Proposition~\ref{prop:report-noise}).

The empirical results are best read as a diminishing-returns scale for case volume within a period. With an apparent balance scale of 37.4 admissions per week on the scale of $R_A$, a state that already reports at the balance scale obtains half of the precision that unlimited case volume could provide, and 90\% requires nine times as many admissions. More than half of the state-weeks in the analysis windows were already above this scale. The limit applies to case volume per period. The error of an estimate pooled over $T$ independent periods falls as $v_R/T$, estimators that shrink towards a prior mean can go below the per-period value (Appendix~\ref{app:shrinkage}), and a power law without an intercept described the observed range of scales about as well as the affine form. The balance scale is therefore a statement about where returns diminish, and it does not imply that error stops falling with more data of every kind. This interpretation is also robust to the source of the intercept. Whatever mixture of transmission, testing and statistic construction it contains, a component that is shared within a period and does not fall with case volume limits the value of further cases per period and raises the value of further periods.

The multiscale decomposition locates this component. About 41\% of the intercept of $R_A$ and 51\% of that of $G$ was common to all states in a calendar week, and resampling weeks moved the estimated intercept 1.5 to 1.9 times as much as resampling states. A national-week component cannot be reduced by scaling up surveillance in a single state, because it enters every state's loss in that week. Pooling states removes the state-specific part, and in these data pooled variance approached the national-week component, about a quarter of single-state variance, with 25 to 33 states. This share is consistent with national influences such as weather, school calendars and holidays, and with national changes in testing or admission practice \citep{dalziel2018urbanization}. Two-way inference is required for the same reason; intervals that resample states only were about two to five times too narrow for the intercept \citep{cameron2011robust}.

Within the limits of an observational comparison, the contrast between reporting regimes is consistent with a benefit of high facility participation. Under mandatory reporting, a median of about 93\% of hospitals reported. Adjustment for reporting completeness changed the intercept of $R_A$ by 2\%, with an upper bound of 5\%, and the change was equivalent to zero within a margin of 17\%. During the voluntary period of 2024, when the median share fell to about 46\%, the interquartile range of log growth was 1.58 times that of the same months in 2023. The interval for this ratio, 0.95 to 2.23, includes one, and the comparison rests on a single off-season. Several checks address the main confounders. The placebo comparison of 2025 with 2023 gave 1.01, matching on reporting scale and on calendar month preserved a ratio above one, and COVID-19 and RSV admissions showed ratios above one in the same period. These results are consistent with an association between low and variable facility participation and more volatile growth statistics, and they do not identify a causal effect of the reporting mandate. The adjustment measures facility participation only. Testing intensity, coding and care seeking within reporting hospitals are not captured. The laboratory proxy shared 23\% of the intercept of $R_A$, but it does not discriminate between transmission and ascertainment, because laboratory growth carries its own noise and attenuates the shared fraction towards zero.

These results suggest a design rule for routine surveillance of the mean reproduction number. Where weekly admissions per jurisdiction already exceed the balance scale, additional resources reduce error more when spent on further periods and on further jurisdictions whose variation is not synchronised (Eq.~(\ref{eq:pooling})) than on further cases per period. The value of complete reporting depends on the reporting volatility a system actually has. In the hypothetical example of Section~\ref{sec:res-theory}, with a relative volatility of $\tau=0.15$, standardising reporting reduced error more than a tenfold increase in case volume per period. The completeness bound observed under mandatory reporting allows $\tau$ up to about 0.04, and at that value a tenfold increase in cases reduced the mean squared error from $15.2\times10^{-3}$ to $11.4\times10^{-3}$, against $14.6\times10^{-3}$ for standardised reporting. In these data, maintaining high facility participation is therefore associated with lower volatility, and its value relative to case volume grows with the volatility that incomplete participation introduces. Seasons are the limiting unit for statements about variation between seasons. With four seasons, no unbiased estimator of a between-season variance can have a relative standard error below 0.82 (Corollary~\ref{cor:season}), and the intercept ratios for alternative denominator weights did not discriminate between mechanisms, because epidemic growth with reporting noise alone already produced $q_1<1<q_2$.

The scope of these conclusions is set by the data and the calibration. The empirical results rest on one surveillance system, four influenza seasons and rule-defined early-growth windows. The window rule was revised once and the renewal-type ratio became the primary statistic during the analysis; Table~\ref{tab:s9} lists each specification with its role and date, and the headline estimates lie within each other's intervals. The inversion is conditional on the generation interval and on the simulated delay and offspring processes, and $R_A$ is a transformation of weekly growth, not a posterior estimate of $R_t$ \citep{cori2013new}. The propositions assume environmental variation that is independent between periods and a fixed number of primary cases; the renewal formulation of Section~\ref{sec:meth-loss} shows that uncertainty in infection pressure adds a further term that does not fall with case volume. The equivalence margins use the sample mean of each statistic, the null ranges for the intercept ratios in rule windows come from the calendar-window simulation, and the bin error bars in Fig.~\ref{fig:floor}B ignore clustering. Three extensions would sharpen these conclusions: renewal estimates with an explicit admission delay and an externally estimated generation interval, hierarchical models with explicit week effects combined with independent measures of testing, and a second surveillance system with mandatory reporting and more seasons.
